\chapter{Reproducibility Details}
\label{ch:appendix}

This appendix collects the hyperparameters, infrastructure, and evaluation-artefact locations referenced throughout Chapters~4--5, in one place, for a reader wanting to reproduce or extend a specific run rather than re-derive its settings from the surrounding narrative.

\section{Training Hyperparameters}
\label{sec:appendix-hyperparams}

\begin{table}[H]
  \caption{Hyperparameters for every trained model reported in Chapter~5.}
  \label{tab:appendix-hyperparams}
  \centering
  \small
  \begin{tabular}{p{2.6cm}p{2.6cm}p{2.6cm}p{2.6cm}p{2.6cm}}
    \toprule
    & Track~1 & Track~2 & DTM estimator & Track~3 ControlNet \\
    \midrule
    Base architecture & CycleGAN (ResNet-9 gen.) & CycleGAN (ResNet-9 gen.) & ControlNet+LoRA & ControlNet (SD1.5) \\
    Training pairs/images & 37{,}054 / 20{,}417 (unpaired) & 1{,}210 (unpaired vs.\ rover) & Gale DTM/ortho patches & 21 real pairs \\
    Epochs / steps & 100 epochs & 25 epochs & 20{,}000 steps & 50 epochs \\
    $\lambda_{\mathrm{cyc}}$ / $\lambda_{\mathrm{idt}}$ & 3.0 / 1.5 & 3.0 / 1.5 & --- & --- \\
    Batch size & 4 & 4 & 4 & 2 \\
    Learning rate & $2{\times}10^{-4}$ & $2{\times}10^{-4}$ & $1{\times}10^{-5}$ & $1{\times}10^{-5}$ \\
    LoRA rank & --- & --- & 16 & --- \\
    Hardware & RTX A4000 & A100 (Eureka2) & A100 (Eureka2) & A100 (Eureka2) \\
    Wall-clock time & \textasciitilde52\,hours & \textasciitilde16\,minutes & --- & --- \\
    \bottomrule
  \end{tabular}
\end{table}

Track~1 and Track~2 share their loss-weight configuration ($\lambda_{\mathrm{cyc}}=3.0$, $\lambda_{\mathrm{idt}}=1.5$), validated by the three-way smoke test reported in Chapter~5, Section~\ref{sec:track1}; the same A100 that completed Track~2's 25 epochs in roughly sixteen minutes took the A4000's RTX A4000 approximately fifty-two hours for Track~1's 100, the empirical basis for this project's general preference for the A4000 on short, queue-free jobs and Eureka2's SLURM cluster on GPU-bound ones, noted throughout Chapter~4.

\section{Evaluation Artefacts}
\label{sec:appendix-artefacts}

Every quantitative result in Chapter~5 is backed by a JSON file recording the metric value, sample sizes, checkpoint identifier, and a timestamp, rather than a number transcribed once and left unverifiable:

\begin{itemize}
  \item \texttt{CycleGAN/Data/eval/} \texttt{fid\_results.json} --- Track~1 FID.
  \item \texttt{CycleGAN/Data/eval/} \texttt{kid\_results\_track1.json}, \texttt{kid\_results\_track2.json} --- Table~\ref{tab:common-metrics}'s KID values.
  \item \texttt{ControlNet/Data/eval/} \texttt{mars\_model\_eval.json} --- Track~3 KID/SSIM/PSNR.
  \item \texttt{ControlNet/Data/eval/} \texttt{dtm\_estimator\_eval.json} --- DTM estimator RMSE.
  \item \texttt{ControlNet/Data/eval/} \texttt{e2e\_pipeline\_metrics.json} --- Table~\ref{tab:e2e-metrics}'s end-to-end SSIM/PSNR/LPIPS values.
\end{itemize}

\section{Infrastructure}
\label{sec:appendix-infra}

Training and evaluation ran across three machines with distinct roles: a local development machine (no CUDA, used for data pipeline development, testing, and disk-heavy fetch steps such as \texttt{prepare\_e2e\_crop.py}), an RTX~A4000 workstation (25\,GB NFS home quota, no SLURM queue, used for short or disk-light GPU jobs), and Surrey's Eureka2 SLURM cluster (A100~80\,GB nodes, used for larger training runs and jobs where queue wait time is acceptable). Checkpoints and generated-image sets were treated as ephemeral and regenerable from a checkpoint plus a fixed evaluation script where disk quota required pruning them, rather than a project asset to be preserved indefinitely; every pruning decision that removed an unevaluated checkpoint is noted at the point in Chapter~5 where its absence matters.
